\chapter{Methodology}
\label{ch:methodology}

\section{Graph Construction}
\label{sec:methodology-graph}

The retrieval systems compared in this thesis operate over an explicit graph built on top of the
IBM Project Debater corpus described in Chapter~\ref{ch:corpus}, rather than over a flat
collection of embedded arguments. At its core, the graph has two node types, policies and
arguments, connected by edges of four types. \textsc{Supports} and \textsc{attacks} edges connect
an argument to the policy it addresses, following directly from the corpus's crowd-annotated
stance label: a PRO argument receives a \textsc{supports} edge to its policy and a CON argument an
\textsc{attacks} edge, each carrying the annotation's stance confidence as an edge weight.
\textsc{Contradicts} and \textsc{equivalent} edges connect pairs of arguments to each other, and,
unlike the stance edges, are not present in the corpus and have to be recovered.
\textsc{Contradicts} edges are the backbone of the graph-aware retrieval system's multi-hop
expansion (Section~\ref{sec:methodology-systems}); \textsc{equivalent} edges are constructed for
the same argument pairs but play a smaller role in this thesis's retrieval systems, and are
retained chiefly as a resource for identifying near-duplicate arguments in the corpus. Above the
policy level, two further node types, MetaPolicy and Topic, group policies hierarchically rather
than argument by argument; this grouping is described next and is used later in the hardened pool
construction (Section~\ref{sec:methodology-hardened-pool}) to identify topically-adjacent sibling
policies for a target policy.

% Suggested figure placement: a small schema diagram here showing all four node types
% (Topic, MetaPolicy, Policy, Argument) and the edge types (BELONGS_TO, SUPPORTS, ATTACKS,
% CONTRADICTS, EQUIVALENT), with stance edges drawn from arguments to their policy and
% CONTRADICTS/EQUIVALENT edges drawn between arguments. This mirrors the "entities and
% relationships" figure in the underlying conference paper and gives the reader a single
% reference point before the NLI recovery pipeline is described below.

\subsection{Policy Grouping: MetaPolicy and Topic}
\label{sec:methodology-metapolicy}

The 71 policies in the corpus are not treated as a flat, unordered list. Two coarser grouping
levels sit above them: a \textsc{belongs\_to} edge connects each Policy node to one of 10
hand-labeled MetaPolicy nodes, and each MetaPolicy node in turn connects via its own
\textsc{belongs\_to} edge to one of 5 Topic nodes. This hierarchy was introduced because several
downstream steps need a principled notion of which policies are topically close to a given
target policy without being the same policy, most directly the sibling-policy selection used to
construct hardened, distractor-injected evaluation pools (Section~\ref{sec:methodology-hardened-pool}):
a distractor argument drawn from a policy in the same MetaPolicy as the target (for example,
drawing a hard negative for ``We should ban cosmetic surgery'' from the sibling policy ``We should
legalize organ trade'', both under Bioethics \& Personal Autonomy) shares vocabulary and framing
with the target policy without answering it, which is precisely what makes it a useful adversarial
test rather than a trivially off-topic one.

Grouping the 71 policies was attempted first as an unsupervised clustering problem, before
falling back to manual assignment, following four steps:

\begin{enumerate}
\item \textbf{Centroid computation.} Each policy was represented by a single centroid vector, the
mean of the \textsc{sentence-transformers/all-mpnet-base-v2} embeddings
\cite{reimers2019sbert,song2020mpnet} of that policy's own PRO and CON arguments, re-normalized to
unit length after averaging since the mean of several unit vectors does not itself have unit norm.
This gives a $71 \times 768$ centroid matrix, one row per policy.

\item \textbf{Dimensionality reduction.} The centroid matrix was reduced to 20 dimensions with
PCA to denoise it before clustering, rather than clustering directly in the full 768-dimensional
space.

\item \textbf{$k$-means sweep.} $k$-means was run on the reduced space for every $k$ from 5 to 15,
each run scored by silhouette coefficient.

\item \textbf{HDBSCAN sweep.} HDBSCAN \cite{campello2013hdbscan} was run on the same reduced space
with its minimum cluster size ranging from 3 to 6, each run scored by silhouette coefficient over
the non-noise points.
\end{enumerate}

Neither method produced a clean, well-separated partition:

\begin{itemize}
\item \textbf{$k$-means} silhouette scores stayed in a narrow and modest band across the entire
swept range, peaking around $k=9$--$10$ at a score of roughly 0.14, well below the 0.5-plus range
conventionally read as indicating real cluster structure.

\item \textbf{HDBSCAN} was, if anything, less decisive: at a minimum cluster size of 6 it labeled
every one of the 71 policies as noise and returned no clusters at all, and even at the smaller
minimum cluster sizes that did return clusters it found only 2--4 of them, with 55--79\% of
policies left unclustered as noise.

\item \textbf{Agreement between the two.} The two methods broadly agreed on the order of magnitude
of a defensible cluster count, roughly single digits to the low teens, but neither converged on a
specific partition. The weak silhouette scores throughout indicate that the 71 policies do not
separate into discrete regions of embedding space so much as blend into one another, which is
unsurprising given that many crowd-authored debate topics share vocabulary or framing across what
are, substantively, different policy questions.
\end{itemize}

Given the absence of a numerically decisive clustering solution, the final MetaPolicy grouping
was assigned manually: the candidate $k$-means and HDBSCAN partitions were used as a reference
point, but each of the 71 policies was ultimately placed into one of 10 hand-labeled MetaPolicy
categories (Table~\ref{tab:metapolicies}) by inspecting the policy texts directly for substantive
thematic overlap, rather than by taking either algorithm's output as final. Scored on the same
reduced embedding space, this manual assignment obtains a silhouette score of 0.09, lower than
$k$-means' best automatic partition; this is expected rather than a failure of the manual step,
since the manual assignment optimizes for a human-legible, substantively coherent category
scheme rather than for geometric separation in embedding space, and the two criteria need not
agree given how weakly separated the policies are to begin with. The 10 MetaPolicies were then
grouped, again manually, under 5 broader Topic categories (Politics \& Governance, Law \& Justice,
Health \& Ethics, Economy \& Society, and Social \& Cultural Issues), giving the three-level
Topic~/~MetaPolicy~/~Policy hierarchy used throughout the rest of this thesis.

\begin{table}[htbp]
\centering
\caption{The 10 manually assigned MetaPolicy categories grouping the corpus's 71 policies.}
\label{tab:metapolicies}
\begin{tabular}{ll}
\toprule
\textbf{ID} & \textbf{MetaPolicy} \\
\midrule
1 & Political Governance \\
2 & Criminal Justice \& Law \\
3 & Civil Liberties \& Rights \\
4 & Bioethics \& Personal Autonomy \\
5 & Economic \& Labor Policy \\
6 & Religion \& Secularism \\
7 & Gender, Family \& Social Norms \\
8 & Education, Media \& Technology \\
9 & Animal Welfare \& Environment \\
10 & Child \& Family Welfare \\
\bottomrule
\end{tabular}
\end{table}

\subsection{Calibrating the \textsc{Contradicts} Threshold}
\label{sec:methodology-edges-calibration}

Recovering \textsc{contradicts} and \textsc{equivalent} edges is framed as a natural language
inference (NLI) problem, following the argumentation literature reviewed in
Section~\ref{sec:background-nlp}: two arguments contradict each other if each entails the
other's negation, and are equivalent if each entails the other directly. An initial attempt at
this idea accepted any candidate pair whose forward contradiction probability alone exceeded a
fixed threshold, and produced a graph in which roughly a third of all possible cross-stance pairs
within a policy were connected by a \textsc{contradicts} edge. Inspecting this result showed the
underlying problem: single-direction contradiction probability is close to 1.0 for nearly any
pair of opposing-stance arguments on the same topic, since an NLI model reads mere stance
opposition as contradiction even when the two arguments do not actually engage the same
underlying claim. This yields a stance-opposition graph rather than a rebuttal graph, and at that
density PageRank centrality and multi-hop traversal would reflect little more than how many
opposing arguments a policy happens to have, not genuine argumentative structure.

Since neither extreme, a near-universally connected stance-opposition graph nor an overly sparse
one that discards genuine rebuttals, is usable, calibration targets a middle band, informally set
at roughly 5--15\% saturation: wide enough to retain a meaningfully connected rebuttal structure
per policy, but far below the 33\% observed under the single-direction rule. Two changes were
available to reach this target: replacing the single-direction acceptance rule with a
\emph{bidirectional} one, requiring the \emph{minimum} of a pair's forward and backward
contradiction probabilities to clear a threshold $\tau_{\text{contra}}$ rather than either
direction alone; and simply raising the cosine floor used to block candidate pairs before NLI is
run at all. It was not obvious in advance which of the two would do the actual work of thinning
the graph, so a dedicated calibration script was run on a small sample of policies, chosen to span
the most-, median-, and least-saturated ends of the original single-direction graph, before
committing to a full 71-policy regeneration. The script reblocks candidate pairs at a deliberately
low cosine floor of 0.40, low enough to expose the full score distribution rather than one already
truncated by the floor being calibrated, scores every candidate pair bidirectionally, and reports
two diagnostics per pair beyond the raw forward and backward contradiction probabilities:

\begin{itemize}
\item \textbf{\emph{contra\_min}}, the minimum of the forward and backward contradiction
probabilities, $\min(P_{\text{fwd}}, P_{\text{bwd}})$: the quantity ultimately thresholded by the
bidirectional acceptance rule.

\item \textbf{Gap}, the absolute difference between the two directions,
$|P_{\text{fwd}} - P_{\text{bwd}}|$: a diagnostic for whether bidirectionality is doing anything
\emph{beyond} what the blocking floor already does. If the gap is consistently close to zero, the
two directions agree with each other almost everywhere, contradiction behaves as a symmetric
property of the pair, and the bidirectional-minimum rule cannot thin the graph beyond what raising
the cosine floor alone would already do. A gap that is often large would instead mean the two
directions frequently disagree, and bidirectional-minimum is doing real, direction-sensitive
filtering of its own.
\end{itemize}

On this corpus, contradiction probability turned out to be a close to binary quantity in each
direction, clustering near 0 or 1 rather than spreading continuously between them, so the gap was
small for the large majority of pairs: the cosine floor, not the bidirectional-minimum rule, is
the dominant lever controlling saturation. The bidirectional rule was kept regardless of this
finding, since it costs nothing to apply and still removes the minority of pairs where the two
directions genuinely disagree.

With this evidence in hand, the calibration script sweeps a small grid of candidate operating
points and reports the resulting saturation at each one:

\begin{itemize}
\item four candidate cosine floors: $0.50$, $0.60$, $0.70$, $0.80$;
\item five candidate \emph{contra\_min} thresholds $\tau_{\text{contra}}$: $0.60$, $0.80$, $0.90$,
$0.95$, $0.99$;
\item for each of the resulting 20 (floor, threshold) pairs, the saturation, edges kept over all
possible cross-stance pairs, in the calibration sample.
\end{itemize}

A floor of 0.70 together with $\tau_{\text{contra}} = 0.90$ was the combination selected, landing
at approximately 11\% saturation on the calibration sample, inside the targeted 5--15\% band.
Table~\ref{tab:edge-calibration} reports the effect of applying this operating point to the full
71-policy graph: the \textsc{contradicts} edge count falls from approximately 1.06 million (33\%
of possible cross-stance pairs, the single-direction rule) to 344,000 (9\%, the bidirectional rule
at this operating point), consistent with the saturation observed on the smaller calibration
sample. This operating point, together with an analogous bidirectional threshold
$\tau_{\text{entail}} = 0.60$ for \textsc{equivalent} edges, is used throughout the four-stage
recovery pipeline described next; paraphrase relations do not exhibit the same saturation problem
that motivated the calibration above, since two arguments rarely entail each other bidirectionally
by accident, but the same bidirectional discipline is applied to them regardless.

\begin{table}[htbp]
\centering
\caption{Effect of the bidirectional-minimum criterion on \textsc{contradicts} edge count,
computed over the full 71-policy graph at a blocking floor of 0.70 and threshold of 0.90.}
\label{tab:edge-calibration}
\begin{tabular}{lrr}
\toprule
\textbf{Acceptance rule} & \textbf{\textsc{Contradicts} edges} & \textbf{Saturation} \\
\midrule
Single-direction threshold & $\approx$1.06M & 33\% \\
Bidirectional minimum & 344{,}000 & 9\% \\
\bottomrule
\end{tabular}
\end{table}

\subsection{Argument-Level Edge Recovery}
\label{sec:methodology-edges}

Applying an NLI model to every pair of arguments within a policy is computationally wasteful,
since the overwhelming majority of pairs are unrelated, so edge recovery is not a single
classification step but a pipeline that narrows the full space of possible argument pairs down to
a small, high-confidence edge set. With the operating point fixed by the calibration in
Section~\ref{sec:methodology-edges-calibration}, a blocking floor of 0.70 and a bidirectional
threshold $\tau_{\text{contra}} = 0.90$ for \textsc{contradicts}, this proceeds in four stages:

\begin{enumerate}
\item \textbf{Blocking.} A cheap embedding-cosine filter, computed over
\textsc{sentence-transformers/all-mpnet-base-v2} embeddings \cite{reimers2019sbert,song2020mpnet},
restricts the pairs that will ever be scored by the far more expensive NLI model. Cross-stance
(PRO/CON) pairs with cosine similarity of at least 0.70 become candidates for
\textsc{contradicts}; within-stance (PRO/PRO or CON/CON) pairs with cosine similarity of at least
0.75 become candidates for \textsc{equivalent}. This step alone discards the large majority of the
$\binom{n}{2}$ possible pairs within a policy, since most arguments on a given policy are neither
close paraphrases nor topically close enough to plausibly stand in either relation.

\item \textbf{Bidirectional NLI classification.} Every surviving candidate pair is scored twice by
the \textsc{cross-encoder/nli-deberta-v3-base} cross-encoder \cite{he2023debertav3}, trained on
SNLI and MultiNLI \cite{bowman2015snli,williams2018mnli}: once with the first argument as premise
and the second as hypothesis, and once with the roles reversed. NLI is a directional task, so a
single scoring direction cannot be assumed to represent the pair as a whole; scoring both
directions is what makes the bidirectional acceptance rule below possible.

\item \textbf{Bidirectional thresholding.} A candidate is accepted as \textsc{contradicts} only if
the \emph{minimum} of its forward and backward contradiction probabilities clears
$\tau_{\text{contra}} = 0.90$, and, symmetrically, as \textsc{equivalent} only if both directions
of entailment probability clear $\tau_{\text{entail}} = 0.60$. An optional third-stage LLM
adjudicator can additionally re-score pairs whose NLI confidence falls in an uncertain middle
band, but this option was left disabled for the edges used throughout this thesis, so that every
reported edge is decided by NLI alone.

\item \textbf{Conflict resolution.} Because \textsc{contradicts} and \textsc{equivalent} are each
symmetric relations, an unordered pair is stored once rather than twice. A small number of pairs
clear both the \textsc{contradicts} and the \textsc{equivalent} threshold simultaneously; for
these, only the higher-confidence verdict is kept, so no argument pair is ever tagged with both
relations at once.
\end{enumerate}

\subsection{PageRank Centrality}
\label{sec:methodology-pagerank}

Once the \textsc{contradicts} edge set is finalized, PageRank \cite{page1999pagerank} is computed
over it to supply the centrality signal used by the graph-aware retrieval system
(Section~\ref{sec:methodology-systems}). PageRank was originally formulated for the directed graph
of hyperlinks between web pages, where the score of a page $p_i$ is defined recursively as

\begin{equation}
\mathrm{PR}(p_i) = \frac{1-d}{N} + d \sum_{p_j \in M(p_i)} \frac{\mathrm{PR}(p_j)}{L(p_j)},
\label{eq:pagerank-classic}
\end{equation}

\noindent where $N$ is the total number of nodes, $M(p_i)$ is the set of nodes linking to $p_i$,
$L(p_j)$ is the out-degree of $p_j$, and $d \in [0,1]$ is a damping factor that mixes the
link-following term against a uniform teleport probability, $\frac{1-d}{N}$, of jumping to any
node at random. In matrix form, writing $\mathbf{v}$ for the uniform teleport vector
($v_i = 1/N$) and $\mathbf{A}$ for the row-normalized transition matrix of the graph, this is the
stationary distribution of

\begin{equation}
\mathbf{PR} = d \, \mathbf{A}^\top \mathbf{PR} + (1-d)\, \mathbf{v},
\label{eq:pagerank-matrix}
\end{equation}

\noindent solved in practice by power iteration rather than direct linear solution, since the
graphs involved are far too large to invert directly.

Two adaptations are made to apply Equations~\ref{eq:pagerank-classic}--\ref{eq:pagerank-matrix} to
the \textsc{contradicts} graph built in Section~\ref{sec:methodology-edges}. First, because
\textsc{contradicts} edges are symmetric by construction (an edge is accepted only if contradiction
holds in both directions, Section~\ref{sec:methodology-edges}), the graph is treated as undirected,
so $M(p_i)$ and $L(p_j)$ in Equation~\ref{eq:pagerank-classic} reduce to a node's neighbors and
degree in the symmetrized \textsc{contradicts} graph rather than to directed in-links and
out-links. Second, alongside a single global PageRank computed with the standard uniform teleport
vector $\mathbf{v}$, a \emph{personalized} PageRank is computed separately for each policy,
replacing $\mathbf{v}$ with a teleport vector $\mathbf{v}_p$ restricted to that policy's own PRO
and CON arguments,

\begin{equation}
v_{p,i} =
\begin{cases}
1 / |A_p| & \text{if } i \in A_p, \\
0 & \text{otherwise},
\end{cases}
\label{eq:pagerank-personalized}
\end{equation}

\noindent where $A_p$ is the set of arguments attached to policy $p$ by a \textsc{supports} or
\textsc{attacks} edge. Restarting the random walk only at policy $p$'s own arguments measures
centrality relative to that single policy's rebuttal structure, rather than to the corpus as a
whole; an argument that is central within a small, tightly contested policy is not penalized for
the fact that a much larger policy elsewhere in the corpus has more \textsc{contradicts} edges in
absolute terms, since the global variant of Equation~\ref{eq:pagerank-classic} would conflate the
two. Both variants use a damping factor of $d=0.85$, the standard choice inherited from the
original formulation, and are computed by power iteration to an L1 residual tolerance of
$10^{-10}$ or a maximum of 200 iterations, whichever comes first. It is the personalized,
per-policy PageRank score, $\mathrm{PR}_p(a) \triangleq \mathrm{PR}(a; \mathbf{v}_p)$, that the
graph-aware retrieval system in Section~\ref{sec:methodology-systems} uses as its relevance
signal.

\section{Retrieval Framework}
\label{sec:methodology-retrieval-framework}

Every retrieval system compared in this thesis, whether it uses the graph built in
Section~\ref{sec:methodology-graph} or not, retrieves stance-aware: given a policy, it returns a
top-$k$ list of PRO arguments and a separate top-$k$ list of CON arguments, so that a single
policy query always yields two independent retrieved lists rather than one undifferentiated one.
Both lists, for every system, are constructed the same way: by selecting $k$ items from a
candidate pool using Maximal Marginal Relevance (MMR) \cite{carbonell1998mmr}, introduced in
Section~\ref{sec:background-rag-howitworks} as the standard mechanism for balancing relevance
against redundancy in a retrieved set. This section fixes the shared MMR mechanics that every
system uses identically; what actually distinguishes one system from another, the candidate pool
each one searches and the relevance signal each one ranks by, is deferred to
Section~\ref{sec:methodology-systems}.

\subsection{Maximal Marginal Relevance}
\label{sec:methodology-mmr}

MMR builds a retrieved list greedily, one item at a time, rather than taking the $k$ most relevant
candidates outright: at each step it adds whichever remaining candidate maximizes a weighted
combination of relevance to the query and dissimilarity to everything already selected. Writing
$R$ for the full candidate pool and $S \subseteq R$ for the set of items already selected
(initially empty), the next item added to $S$ is

\begin{equation}
\operatorname*{arg\,max}_{d_i \in R \setminus S}
\left[ \lambda \cdot \mathrm{rel}(d_i) - (1-\lambda) \max_{d_j \in S} \mathrm{sim}(d_i, d_j) \right],
\label{eq:mmr}
\end{equation}

\noindent where $\mathrm{rel}(d_i)$ is candidate $d_i$'s relevance score, defined per system in
Section~\ref{sec:methodology-systems}; $\mathrm{sim}(d_i, d_j)$ is embedding cosine similarity
between two candidates, computed over the same
\textsc{sentence-transformers/all-mpnet-base-v2} embeddings \cite{reimers2019sbert,song2020mpnet}
used throughout this thesis; and $\max_{d_j \in S} \mathrm{sim}(d_i, d_j)$ is therefore the
similarity between $d_i$ and whichever already-selected item is closest to it, the redundancy
penalty. The weight $\lambda \in [0,1]$ controls the relevance/diversity tradeoff directly:
$\lambda = 1$ reduces Equation~\ref{eq:mmr} to plain top-$k$ ranking by relevance alone, with no
diversification, while $\lambda = 0$ ignores relevance entirely and greedily maximizes spread
against what has already been picked. Every system and every experiment in this thesis fixes
$\lambda = 0.5$, giving relevance and diversity equal weight; because this value never changes
across systems, any diversity difference observed in the results (Chapter~\ref{ch:results}) is
attributable to differences in the candidate pool $R$ or the relevance signal $\mathrm{rel}(\cdot)$
between systems, not to one system having been handed a more favorable relevance/diversity
tradeoff than another.

The first item MMR selects is a useful special case for making Equation~\ref{eq:mmr} concrete: with
$S = \emptyset$, the redundancy term vanishes, and the first pick reduces to
$\operatorname*{arg\,max}_{d_i \in R} \mathrm{rel}(d_i)$, plain top-1 relevance ranking. Every
subsequent pick then balances that same relevance signal against dissimilarity to whatever has
been selected so far, so a system's relevance signal alone determines where retrieval starts, and
the interaction between that signal and the pool's own similarity structure determines everywhere
it goes after that.

Two design choices are deliberately left open by Equation~\ref{eq:mmr}, and it is exactly these
two that distinguish one retrieval system from another in this thesis:

\begin{itemize}
\item how the candidate pool $R$ is generated, that is, which arguments are even eligible to be
retrieved before MMR is run at all; and

\item what relevance signal $\mathrm{rel}(\cdot)$ ranks the pool by.
\end{itemize}

Holding the corpus, the embedding model, the MMR mechanism, and $k$ fixed across every system, and
varying only these two choices, is what makes it possible to attribute an observed difference in
retrieval quality to a specific design decision, candidate generation or relevance ranking, rather
than to a confounded bundle of implementation differences. Section~\ref{sec:methodology-systems}
defines each system compared in this thesis exactly by how it answers these two questions.

\section{Retrieval Systems Compared}
\label{sec:methodology-systems}

Four retrieval systems are compared throughout this thesis, all built on the shared MMR mechanism
of Section~\ref{sec:methodology-retrieval-framework} and differing in exactly one component at a
time, the candidate pool $R$, the relevance signal $\mathrm{rel}(\cdot)$, or the diversification
step itself, holding the corpus, the embedding model, and $k$ fixed across all four. This
single-variable design is what licenses attributing an observed difference between two systems to
the one component that actually differs between them, rather than to a confounded bundle of
implementation choices.

\subsection{StandardRAG: the Embedding-Only Baseline}
\label{sec:methodology-standardrag}

StandardRAG is the embedding-only baseline against which the graph-aware system is compared, and
uses no graph signal at any stage.

\begin{itemize}
\item \textbf{Candidate pool.} The pool $R$ is the policy's own native argument set, every
argument already carrying a \textsc{supports} or \textsc{attacks} edge to the policy
(Section~\ref{sec:methodology-graph}), capped at 200 per stance. No query is embedded and no
corpus-wide search is performed to build this pool: construction is a direct lookup of the
policy's own tagged arguments, truncated to 200 in the order they are stored rather than ranked
by similarity to anything. Cosine similarity enters only at the relevance-signal stage below, not
at candidate-generation time.

\item \textbf{Relevance signal.} Once the pool is fixed, candidates are ranked by cosine
similarity, over \textsc{sentence-transformers/all-mpnet-base-v2} embeddings
\cite{reimers2019sbert,song2020mpnet}, to $\bar{d}$, the centroid of that same pool:
$\mathrm{rel}(d_i) = \mathrm{sim}(d_i, \bar{d})$. This is a form of pseudo-relevance feedback,
ranking candidates by centrality to an already-assembled candidate set rather than by similarity
to an explicit query, and it is the only point at which embedding similarity plays any role in
StandardRAG's pipeline.
\end{itemize}

Because StandardRAG's candidate pool is built directly from the \textsc{supports}/\textsc{attacks}
edges already present in the graph (Section~\ref{sec:methodology-graph}), rather than from a
separate corpus-wide embedding search, it is graph-\emph{adjacent} in its pool construction even
though it consults no \textsc{contradicts} structure and computes no centrality score; ``no graph
signal'' refers specifically to the relevance signal and to the absence of any multi-hop
expansion, not to the origin of the pool itself. This matters for how StandardRAG compares to
PR-GraphRAG below: under the native pool, the two systems' candidate pools share the same
direct-attachment seed, and differ only in whether that seed is expanded by one
\textsc{contradicts} hop; the axis that separates them most consequentially under this condition
is therefore the relevance signal, cosine-to-centroid versus personalized PageRank, with the pool
itself diverging only at the margin, for policies whose native argument count for a stance exceeds
the 200 cap.

\subsection{PR-GraphRAG: the Graph-Aware System}
\label{sec:methodology-prgraphrag}

PR-GraphRAG is this thesis's proposed system, and is the only one of the four that consults the
\textsc{contradicts} structure of the graph built in Section~\ref{sec:methodology-graph}. Its
candidate pool starts from the identical direct-attachment seed as StandardRAG's
(Section~\ref{sec:methodology-standardrag}), so, under the native pool, the two systems are not
distinguished by one having access to a wider or more permissive net over the corpus than the
other; both are structurally confined to a single policy's own tagged arguments at the outset.
What PR-GraphRAG adds is a further expansion of that same seed along \textsc{contradicts} edges,
and a relevance signal, personalized PageRank, that reflects the rebuttal structure connecting
those arguments to one another rather than their raw similarity to a pool centroid. The
motivation this is built to address, an embedding-only signal conflating \emph{an argument that
addresses this policy} with \emph{an argument that merely uses this policy's vocabulary}, becomes
directly relevant once the hardened, distractor-injected pools of
Section~\ref{sec:methodology-hardened-pool} deliberately introduce arguments from topically
adjacent sibling policies (for instance, several policies on cosmetic surgery, organ trade, and
human cloning all sharing a Bioethics \& Personal Autonomy MetaPolicy,
Section~\ref{sec:methodology-metapolicy}) into a policy's own candidate pool: at that point, an
embedding-only relevance signal has no basis for distinguishing a genuine argument from a
vocabulary-sharing distractor, whereas an argument's position, or absence, in the
\textsc{contradicts} structure surrounding the policy's genuine arguments does.

\begin{itemize}
\item \textbf{Candidate pool.} The pool is seeded by the arguments holding a direct
\textsc{supports} or \textsc{attacks} edge to the target policy, then expanded one hop along
\textsc{contradicts} edges: a CON argument that contradicts a directly-attached PRO argument is
added to the PRO candidate pool, and symmetrically, a PRO argument that contradicts a
directly-attached CON argument is added to the CON pool. Because \textsc{contradicts} edges are
recovered only within a policy (Section~\ref{sec:methodology-edges}), this hop cannot itself reach
outside the target policy on the native pool; its effect there is to recover same-policy arguments
that the 200-item truncation of the direct seed may have dropped, not to widen the search to other
policies. The hop's exclusionary role, arguments with no path into the target policy's own
\textsc{contradicts} neighborhood are never added at all, becomes consequential once the pool
itself is allowed to contain material from outside the policy, as it deliberately does under the
hardened-pool construction of Section~\ref{sec:methodology-hardened-pool}.

\item \textbf{Relevance signal.} Candidates are ranked by their personalized, per-policy PageRank
score (Section~\ref{sec:methodology-pagerank}), min--max normalized within the pool, so
$\mathrm{rel}(d_i) = \widehat{\mathrm{PR}}_p(d_i)$.
\end{itemize}

Under Equation~\ref{eq:mmr}, the first item PR-GraphRAG selects is therefore the most central
argument in the policy's rebuttal structure, the argument the largest weight of other arguments
engages with, rather than the argument closest to a pool centroid.

\subsection{RAG-no-MMR: a Diversification Ablation}
\label{sec:methodology-ragnommr}

RAG-no-MMR isolates the contribution of MMR's diversification step in isolation from every other
design choice. It uses the identical candidate pool and the identical relevance signal as
StandardRAG (Section~\ref{sec:methodology-standardrag}), the policy's native argument set capped
at 200, ranked by cosine similarity to the pool centroid, but selects its final top-$k$ list by
plain ranking on $\mathrm{rel}(\cdot)$ alone rather than by running MMR (Equation~\ref{eq:mmr})
over the pool. Formally, this is the $\lambda = 1$ boundary
case of Equation~\ref{eq:mmr} applied specifically to StandardRAG's pool and relevance signal,
rather than a separately-tuned system. Because pool and relevance signal are held fixed and only
the diversification step is removed, any difference between StandardRAG and RAG-no-MMR in
Chapter~\ref{ch:results} isolates exactly what MMR's redundancy penalty buys, or costs, on this
corpus.

\subsection{Random: a Sanity Floor}
\label{sec:methodology-random}

Random selects its top-$k$ list uniformly at random from the same native, policy-restricted pool
that StandardRAG and RAG-no-MMR draw from, using no relevance signal and no diversification step
at all. It serves two purposes in this thesis's evaluation: a sanity floor that any system using a
genuine relevance signal ought to outperform, and, as Chapter~\ref{ch:results} returns to, a
diagnostic for whether a given evaluation pool is discriminative enough to separate a real
retrieval system from one making no retrieval decision whatsoever.

\subsection{Summary}
\label{sec:methodology-systems-summary}

Table~\ref{tab:systems-summary} summarizes all four systems along the two axes fixed in
Section~\ref{sec:methodology-retrieval-framework}, candidate pool and relevance signal, plus
whether MMR diversification is applied. StandardRAG and RAG-no-MMR share an identical pool and
relevance signal, differing only in the last column, isolating MMR's contribution alone.
StandardRAG and PR-GraphRAG share the same direct-attachment seed for their pool (both start from
a policy's own \textsc{supports}/\textsc{attacks} arguments) and differ in whether that seed is
expanded by a \textsc{contradicts} hop and in which relevance signal ranks it, so that any
observed difference between the two is attributable to the graph-derived expansion and ranking
specifically, not to a pool built by an unrelated mechanism. This table is, in effect, a compact
statement of the single-variable design that Chapter~\ref{ch:results} exploits to attribute
observed differences to specific components rather than to confounded implementation choices.

\begin{table}[htbp]
\centering
\small
\caption{The four retrieval systems compared in this thesis, defined along the two axes MMR
leaves open (candidate pool, relevance signal) plus whether diversification is applied.}
\label{tab:systems-summary}
\begin{tabular}{@{}p{2.1cm}p{3.6cm}p{2.9cm}p{1.7cm}@{}}
\toprule
\textbf{System} & \textbf{Candidate pool $R$} & \textbf{Relevance signal $\mathrm{rel}(\cdot)$} & \textbf{MMR} \\
\midrule
StandardRAG & Native pool (\textsc{supports}/\textsc{attacks} seed), capped at 200 & Cosine to pool centroid & Yes ($\lambda=0.5$) \\
PR-GraphRAG & Same seed $+$ 1-hop \textsc{contradicts} expansion & Personalized PageRank & Yes ($\lambda=0.5$) \\
RAG-no-MMR & Native pool (\textsc{supports}/\textsc{attacks} seed), capped at 200 & Cosine to pool centroid & No (plain top-$k$) \\
Random & Native pool (\textsc{supports}/\textsc{attacks} seed), capped at 200 & None (uniform) & No (random sample) \\
\bottomrule
\end{tabular}
\end{table}

\section{Use Cases}
\label{sec:methodology-usecases}

Retrieval quality is evaluated through two downstream use cases rather than one, chosen to
reflect two genuinely different ways a moderator or citizen consumes retrieved arguments: scanning
a short, high-precision list in full, and reading a longer synthesized account grounded in a wider
set of retrieved material. The two use cases share every retrieval system, the corpus, and the
candidate-pool construction described in Sections~\ref{sec:methodology-graph}--\ref{sec:methodology-systems};
they differ only in $k$, in whether the retrieved set is read directly or first summarized, and in
the downstream metrics each one is scored on (Chapter~\ref{ch:evaluation}).

\subsection{UC1: Conflict-Map Retrieval}
\label{sec:methodology-uc1}

Given a policy, UC1 asks a retrieval system to return a short, stance-balanced set of PRO and CON
arguments that a moderator can scan directly to see where a debate's genuine points of
disagreement lie, without having to read the corpus itself. This is a \emph{precision-first} task:
because the retrieved set is read in full rather than filtered further by a human or a downstream
model, every irrelevant or redundant item carries a direct cost, and diversity, avoiding
near-duplicate arguments that restate the same point, matters as much as raw relevance. UC1 is
evaluated over the candidate pool of 200 fixed in Section~\ref{sec:methodology-systems}, and
reported primarily at $k=5$, the operating point at which a human moderator would realistically
read every retrieved item rather than skimming a long list; Chapter~\ref{ch:results} additionally
reports a sweep across pool sizes and $k$ values to check whether conclusions drawn at this single
operating point generalize.

\subsection{UC2: Debate Summarization}
\label{sec:methodology-uc2}

Given a policy, UC2 asks a system to generate a short natural-language summary of the debate,
covering the main PRO and CON positions in balance, grounded in that system's retrieved arguments
rather than in the corpus at large. Unlike UC1, this is a \emph{recall-oriented},
generation-downstream task: retrieval quality matters here not because a person reads the
retrieved list directly, but because it bounds what the summarizer can faithfully say about the
debate. A summarizer cannot surface a position it was never given, so UC2 evaluates whether a
retrieval system's choices survive an additional generation step rather than being read off the
retrieved list itself. UC2 is evaluated at $k=15$, a wider retrieval budget than UC1's, appropriate
for feeding a summarization prompt with enough material to synthesize from rather than for direct
human reading.

Summaries are generated with a structured prompt, issued to Qwen3:8b at temperature 0, that
requires the model to return a JSON object with exactly five fields, rather than free-form prose:
the strongest PRO argument, the strongest CON argument, the key tension between them stated as
specific contradicting claims rather than a restatement of the policy topic, any missing
perspectives, and an overall assessment of how balanced the retrieved debate is. Constraining the
summarizer's output to this fixed schema keeps the four downstream UC2 metrics
(Chapter~\ref{ch:evaluation}) computable consistently across systems and policies, since every
summary decomposes into the same five comparable fields regardless of which retrieval system fed
it.

\section{Policy Selection}
\label{sec:methodology-policy-selection}

Every result reported in this thesis is computed over a fixed sample of 30 policies out of the
corpus's 71, selected once and reused unchanged across every system, pool condition, and metric,
so that no later comparison is confounded by different systems having been evaluated on different
policies.

The obvious selection rule, simply taking the 30 policies with the most arguments, was rejected
because it would bias the sample toward corpus outliers and foreclose any later analysis of how
retrieval performance varies with graph density: Section~\ref{sec:corpus-ibm-stats} already
established a roughly 1.7-fold spread in argument count across the corpus's 71 policies (331 to
554), and a large-policies-only sample would systematically discard the sparse end of that range,
the exact regime in which a graph-aware system's one-hop expansion has the least material to work
with. A tercile-stratified sample was used instead:

\begin{enumerate}
\item The 71 policies are sorted by total argument count and split into three bins of equal size,
sparse, medium, and dense, 23, 23, and 25 policies respectively (71 does not divide evenly by
three).

\item 10 policies are drawn uniformly at random from each bin, with a fixed random seed
($\text{seed}=42$, the same seed used throughout this thesis's other randomized steps, for
consistency) for reproducibility, giving $10 \times 3 = 30$ policies in total.
\end{enumerate}

This produces three evenly-sized groups spanning the corpus's full density range by construction:

\begin{itemize}
\item \textbf{Dense tercile} (432--554 arguments per policy): 10 policies drawn from the 25
policies with the most native argument volume.

\item \textbf{Medium tercile} (405--432 arguments per policy): 10 policies drawn from the 23
policies in the middle of the distribution.

\item \textbf{Sparse tercile} (331--404 arguments per policy): 10 policies drawn from the 23
policies with the least native argument volume, the regime in which a graph-aware system's
one-hop \textsc{contradicts} expansion has the least material to work with.
\end{itemize}

Because the sample is balanced across these three groups by design, any relationship later
observed between graph density and retrieval quality can be attributed to density itself rather
than to an artifact of oversampling one density regime over another; this is what licenses the
density-versus-performance analysis in Chapter~\ref{ch:discussion}.

\section{Hardened Pool Construction}
\label{sec:methodology-hardened-pool}

\subsection{Why the Native Pool Is Not a Sufficient Test}
\label{sec:methodology-hardened-motivation}

Every candidate pool discussed so far, whether used by StandardRAG, PR-GraphRAG, RAG-no-MMR, or
Random, is drawn from a policy's own native argument set (Section~\ref{sec:methodology-systems}).
Inspecting the ground-truth relevance judgments over this native pool (the same judgments
described in Section~\ref{sec:methodology-groundtruth} below) reveals that it is, by construction,
largely relevant already: averaged over the stratified 30-policy sample
(Section~\ref{sec:methodology-policy-selection}), 57.7\% of judged arguments in a native pool
receive a relevance grade of 2 or higher, only 6.8\% receive a grade of 0, and this relevance
density varies almost not at all from one policy's pool to another, a standard deviation of just
0.071 across the 30 sampled policies. The problem this creates is not specific to any one policy,
it is systemic: because most arguments in a native pool are at least somewhat relevant simply by
virtue of having been authored for that policy, a ranking metric has little room to separate a
genuinely good retrieval system from a poor one, or, in the most extreme case, from Random
(Section~\ref{sec:methodology-random}) itself. A relevance signal that does essentially nothing
can still score respectably on a pool where most of what it could possibly return already clears
the relevance bar.

This motivates deliberately constructing a \emph{harder} evaluation pool rather than relying on
the native one, following the broader argument-retrieval literature's recognition that the
specific distractors present in a retrieval context materially affect how discriminating an
evaluation is, and that a retrieval evaluation should therefore construct adversarial candidate
pools rather than rely on native, largely-relevant ones.\footnote{This methodological point
follows Cuconasu et al.'s work on the effect of distractors in retrieval-augmented generation
contexts (SIGIR 2024); a full citation for this source still needs to be added to the thesis
bibliography.} Concretely, the aim is to inject arguments that are topically adjacent to a
policy, close enough in subject matter that an embedding-only signal cannot trivially reject them,
but that do not actually address the policy in question, so that a system's relevance signal is
genuinely tested on its ability to discriminate rather than credited for a pool that is easy by
construction.

\subsection{Constructing the Distractor Pool}
\label{sec:methodology-hardened-pipeline}

For a given target policy $P$ and stance $s \in \{\text{PRO}, \text{CON}\}$, distractors are
constructed in four steps:

\begin{enumerate}
\item \textbf{Identify sibling policies.} Siblings of $P$ are the other policies sharing $P$'s
MetaPolicy (Section~\ref{sec:methodology-metapolicy}), for instance ``We should legalize organ
trade'' and ``We should ban human cloning'' as siblings of ``We should ban cosmetic surgery''
under Bioethics \& Personal Autonomy. Where a policy's MetaPolicy grouping is unavailable, a
fallback identifies siblings instead as the four policies whose centroid
(Section~\ref{sec:methodology-metapolicy}) is nearest by cosine similarity to $P$'s own centroid.

\item \textbf{Gather candidates.} All same-stance arguments belonging to $P$'s sibling policies are
pooled as candidate distractors: sibling PRO arguments feed $P$'s PRO distractor pool, sibling CON
arguments feed $P$'s CON distractor pool. Any candidate duplicating, by text, an argument $P$
already owns is discarded.

\item \textbf{Rank and select.} Candidates are ranked nearest-first by cosine similarity to $P$'s
own centroid, the hardest and most topically confusable distractors, and the top $D$ are kept for
each injection level $D$ (Section~\ref{sec:methodology-hardened-levels}). A random-selection
variant, drawing uniformly from the sibling candidate pool instead of nearest-first, is also
available as a system-agnostic control, to check that any robustness result is not an artifact of
the specific nearest-first selection strategy.

\item \textbf{Naturalize into the target policy.} Each selected distractor is not moved but
\emph{copied}: a fresh argument record is created with a new identifier, re-stamped to belong to
$P$ (\texttt{policy}~$=P$) and to carry stance $s$, and flagged \texttt{is\_distractor}~$=$~true,
with the original policy and argument id retained as provenance fields. Re-stamping is what makes
the distractor copy visible to every retrieval system on equal footing: every candidate-pool
lookup in Section~\ref{sec:methodology-systems} keys off the target policy's own tagged arguments,
so an argument that still carried its original policy label would never enter any system's
candidate pool at all, defeating the purpose of injecting it.
\end{enumerate}

An early version of this pipeline appended distractor copies to the end of each policy's native
argument list, after the (up to 200) native arguments already occupying the first positions. Since
every candidate-pool lookup in Section~\ref{sec:methodology-systems} truncates that list to 200,
the appended distractors, sitting at position 200 and beyond, were silently sliced off before any
system ever saw them: a sanity check that had the Random baseline (Section~\ref{sec:methodology-random})
retrieve zero distractors from a pool that was supposed to be roughly one-third distractors by
construction, a result close enough to statistically impossible that it exposed the bug rather
than confirming robustness. The fix interleaves the native and distractor identifiers with a
deterministic, per-(policy, stance, level) seeded shuffle before any truncation happens, so a
downstream $[:200]$ slice draws a representative native/distractor mixture at the intended
injected rate. This interleaving step is applied only when distractors have actually been
injected, so the zero-distractor control condition below remains byte-identical to the native
pool.

\subsection{Injection Levels}
\label{sec:methodology-hardened-levels}

Distractors are injected at five levels per stance, $D \in \{0, 50, 100, 200, \text{full}\}$, where
\emph{full} takes every available sibling candidate rather than capping at a fixed count. $D=0$ is
a byte-identical control, reproducing the native pool exactly, since no injection and no
interleaving shuffle occur at this level. The results reported in Chapter~\ref{ch:results} focus
primarily on $D=100$ (denoted L100), chosen as a middle setting large enough to meaningfully
stress-test a pool of up to 200 native arguments without the distractor share overwhelming it
entirely, with the full $D$ sweep reported separately to check that conclusions drawn at L100
generalize across injection levels.

\subsection{Four Edge-Connectivity Scenarios}
\label{sec:methodology-hardened-scenarios}

Injecting distractor arguments raises a further question that is easy to overlook: once a
distractor copy exists inside the target policy's argument set, should it participate in the
\textsc{contradicts}/\textsc{equivalent} edge-recovery pipeline of Section~\ref{sec:methodology-edges}
the same way a native argument does? Answering this only one way would conflate two separate
questions, whether graph-aware retrieval resists distractors at all, and specifically whether that
resistance depends on distractors being graph-\emph{isolated}. Four scenarios are therefore
compared below, distinguished by whether, and how, distractor arguments are connected into the
\textsc{contradicts} graph.

The first scenario, \textbf{no distractors}, is the native pool of
Section~\ref{sec:methodology-hardened-motivation}, unmodified.

The second, \textbf{isolated distractors}, is the primary hardened condition used throughout
Chapter~\ref{ch:results}: distractor copies are injected exactly as described in
Section~\ref{sec:methodology-hardened-pipeline}, but the edge-recovery pipeline of
Section~\ref{sec:methodology-edges} is never re-run over them, so they carry no
\textsc{contradicts} or \textsc{equivalent} edge to anything, native argument or fellow distractor
alike. Under this scenario, a distractor can still enter a system's candidate pool directly, since
candidate-pool lookups key off the \textsc{supports}/\textsc{attacks} re-stamping from
Section~\ref{sec:methodology-hardened-pipeline} rather than off \textsc{contradicts} connectivity,
but it can never be reached by PR-GraphRAG's one-hop expansion, and its personalized PageRank
score is that of an isolated node.

The third and fourth scenarios both re-run the full three-stage edge-recovery pipeline of
Section~\ref{sec:methodology-edges} over the merged native-plus-distractor argument set, at the
same calibrated thresholds (Section~\ref{sec:methodology-edges-calibration}), but differ in one
blocking-stage rule. In \textbf{distractors with edges}, every pair is eligible for
\textsc{contradicts}/\textsc{equivalent} scoring exactly as in the native pipeline, so two
distractor copies, both drawn from sibling policies that may genuinely discuss related claims, can
end up connected to each other as readily as two native arguments can. In \textbf{distractors, no
distractor--distractor edges}, the blocking stage additionally skips any candidate pair where both
endpoints are distractor copies, before NLI is ever run on that pair, while leaving
native--distractor pairs to be scored normally; native--native pairs are unaffected in either
scenario, so the two are directly comparable and differ only in whether distractors can connect to
each other.

Table~\ref{tab:hardened-scenarios} summarizes the resulting connectivity of all four scenarios;
Table~\ref{tab:distractor-edge-breakdown} just below it reports the actual edge counts this
produced at L100.

\begin{table}[htbp]
\centering
\caption{The four pool-construction scenarios compared in this thesis, by whether injected
distractor arguments participate in the \textsc{contradicts} graph.}
\label{tab:hardened-scenarios}
\begin{tabular}{lcc}
\toprule
\textbf{Scenario} & \textbf{Distractor--distractor edges} & \textbf{Native--distractor edges} \\
\midrule
No distractors (native) & N/A & N/A \\
Isolated distractors & None & None \\
Distractors with edges & Unrestricted & Yes \\
No distractor--distractor edges & None & Yes \\
\bottomrule
\end{tabular}
\end{table}

Run over the full merged native-plus-distractor argument set at L100, the unrestricted pipeline
(\textbf{distractors with edges}) accepts 385,872 \textsc{contradicts} edges in total. Re-running
the identical pipeline with the additional both-endpoints-distractor exclusion
(\textbf{distractors, no distractor--distractor edges}) accepts 345,109 \textsc{contradicts} edges
and 13,302 \textsc{equivalent} edges, with zero pairs triggering the Stage 4 conflict-resolution
rule (Section~\ref{sec:methodology-edges}) in this run. The 385,872-edge unrestricted total
decomposes into 344,556 native--native edges, matching the native-corpus count already established
in Table~\ref{tab:edge-calibration}; 553 native--distractor edges; and 40,763
distractor--distractor edges, the difference removed by the no-distractor--distractor-edges rule.

\begin{table}[htbp]
\centering
\caption{Decomposition of the 385,872 \textsc{contradicts} edges generated over the merged
native-plus-distractor argument set at L100, by pair type.}
\label{tab:distractor-edge-breakdown}
\begin{tabular}{lr}
\toprule
\textbf{Pair type} & \textbf{\textsc{Contradicts} edges} \\
\midrule
Native--native & 344,556 \\
Native--distractor & 553 \\
Distractor--distractor & 40,763 \\
\midrule
Total (unrestricted) & 385,872 \\
\bottomrule
\end{tabular}
\end{table}

This breakdown is itself informative beyond justifying the four-scenario design: distractor
arguments contradict \emph{one another} at nearly two orders of magnitude the rate at which they
contradict the target policy's own genuine arguments, 40,763 distractor--distractor edges against
only 553 native--distractor edges. Sibling-policy distractors, drawn from topically adjacent
policies under the same MetaPolicy (Section~\ref{sec:methodology-metapolicy}), evidently argue
with each other far more often than they engage the specific arguments already native to the
target policy. This matters for interpreting the difference between scenarios 3 and 4 in
Chapter~\ref{ch:results}: whatever connectivity the unrestricted scenario adds over the
no-distractor--distractor-edges scenario is overwhelmingly internal to the injected distractor set
itself, not a new pathway connecting distractors into the target policy's own argumentative
structure.

Of the 72,349 distractor--distractor candidate pairs that survived blocking but were excluded from
NLI scoring entirely under the no-distractor--distractor-edges rule, roughly 56\% (40,763 of
72,349) would have cleared the bidirectional \textsc{contradicts} threshold had they been scored,
confirming that sibling-policy distractors frequently stand in a genuine contradiction relation to
\emph{each other}, exactly the connectivity that scenario 4 exists to exclude and that scenario 3
exists to test the consequences of. As a validation step before this edge set was used to
recompute PageRank, the edge-count refresh script was run with the expected \textsc{contradicts}
total passed explicitly as a guard (\texttt{--expected-contradicts 345109}), so any mismatch
between the freshly generated edge file and the count reported by Stage 3 would halt the pipeline
rather than silently propagate into the centrality scores PR-GraphRAG depends on.

Comparing all four scenarios makes it possible to attribute any observed robustness specifically:
whether graph-aware retrieval resists distractors at all (scenario 2 versus scenario 1), whether
that resistance is contingent on distractors remaining graph-isolated (scenario 3 versus scenario
2), and, if resistance degrades under scenario 3, whether the cause is distractor--distractor
connectivity specifically rather than distractor connectivity in general (scenario 4 as the
isolating comparison against scenario 3).

\section{Ground Truth Generation}
\label{sec:methodology-groundtruth}

Every metric reported in Chapters~\ref{ch:evaluation}--\ref{ch:results} that depends on a notion
of relevance, most directly UC1's nDCG, is only as trustworthy as the relevance judgments it is
computed against. This section describes how those judgments were produced: an exhaustive,
dual-annotator grading protocol applied to every candidate pool, native and hardened alike.

\subsection{Grading Protocol}
\label{sec:methodology-groundtruth-protocol}

Each (policy, stance) candidate pool is graded on a four-point ordinal relevance scale, 0
(irrelevant) to 3 (essential), by two large language models from different model families, Llama
3.3 70B \cite{grattafiori2024llama3} as the primary annotator and Gemma 31B
\cite{gemma2026gemma4} as the secondary. Two design choices in this protocol depart from a more
convenient alternative, and both are deliberate rather than incidental:

\begin{itemize}
\item \textbf{The full pool is judged, not a shortlist.} Every candidate in a policy's pool, up to
the same 200-argument cap used by the retrieval systems themselves
(Section~\ref{sec:methodology-systems}), receives a grade, rather than only the first $n$
arguments encountered or only the union of what the compared systems happen to retrieve. Judging
a shortlist instead would silently make relevance a positional artifact, an argument sitting past
the shortlist boundary can never be credited as relevant no matter how strong it is, and would
under-judge exactly the tail region a system's own retrieved list might reach into. Grading the
full pool removes that bias, at the cost of a long tail of grade-0 judgments and a real risk of
annotator drift across the several batches a 200-argument pool requires; the per-policy grade
distribution is retained in the full grading report specifically so that such drift remains
visible rather than silently absorbed into an aggregate statistic.

\item \textbf{Two annotators, not one.} A single LLM annotator's grades cannot be distinguished
from that model's idiosyncratic biases. Using two annotators from different model families, and
reporting their agreement (Section~\ref{sec:methodology-groundtruth-agreement}), makes any such
bias visible rather than silently baked into the ground truth. This follows established results
showing that LLM assessors can approach trained human judges' agreement on graded relevance
judgments \cite{thomas2024llmjudge,upadhyay2024umbrela,upadhyay2025largescale}, while a
single-annotator design carries a documented risk of systematic, model-specific judging bias
\cite{panickssery2024selfpref}.
\end{itemize}

\subsection{Consensus Resolution}
\label{sec:methodology-groundtruth-consensus}

With two annotators graded independently, a single ground-truth grade per argument is derived by a
documented consensus rule rather than by an arbitrary tie-break. The rule adopted is
\emph{minimum}: when the two annotators disagree, the lower of the two grades is kept, following
the disagreement heuristic used by the TripJudge collection \cite{althammer2022tripjudge}, on the
reasoning that relevance the two annotators could not agree on should not be treated as though it
were confidently established. An argument's relevance grade, the primary annotator's grade, and
the secondary annotator's grade are all retained in the stored ground truth, so the consensus
resolution is transparent and reversible rather than discarding the individual judgments once
resolved. A binary relevant set, used wherever a metric needs a relevant/non-relevant cut rather
than a graded score, is defined as grade $\geq 2$.

\subsection{Inter-Annotator Agreement}
\label{sec:methodology-groundtruth-agreement}

Agreement between the two annotators is reported at two granularities, following FiRA
\cite{hofstatter2020fira} and TripJudge's \cite{althammer2022tripjudge} observation that
agreement on the full ordinal scale and agreement on the binary relevant/non-relevant cut can
diverge substantially, and that the binary cut is the granularity that actually determines the
relevant set most metrics are computed against:

\begin{itemize}
\item \textbf{Ordinal agreement}, quadratic-weighted Cohen's kappa \cite{cohen1968kappa} over the
full 0--3 grade scale, penalizing a disagreement of, for instance, grade 0 versus grade 3 far more
than a disagreement of grade 1 versus grade 2. Computed per policy and averaged over the
stratified 30-policy sample (Section~\ref{sec:methodology-policy-selection}), this reaches a mean
of 0.584 (std 0.082, range 0.439--0.750 across policies), which the conventional bands of Landis
and Koch \cite{landis1977kappa} place in the \emph{moderate} agreement range on average, with
individual policies ranging as high as the \emph{substantial} band.

\item \textbf{Binary agreement}, unweighted Cohen's kappa over the collapsed relevant ($\geq 2$)
versus non-relevant ($< 2$) cut, reported separately since a metric like nDCG's underlying
relevant set is defined at this granularity rather than the full ordinal one. This reaches a mean
of 0.528 (std 0.081, range 0.348--0.678), somewhat lower and more variable than the ordinal
figure, consistent with a binary cut discarding some of the ordinal scale's partial-credit
structure.
\end{itemize}

Applying the annotation pipeline's own operational thresholds (kappa $\geq 0.6$ = HIGH, $0.4$--$0.6$
= MEDIUM, $<0.4$ = LOW) to the ordinal figure, 13 of the 30 sampled policies fall in the HIGH band
and the remaining 17 in MEDIUM, with none falling into LOW. For 3 of the 30 policies (``We should
subsidize vocational education,'' ``We should close Guantanamo Bay detention camp,'' and ``We
should abandon the use of school uniform''), the secondary annotator's grades were unavailable for
that policy, so no kappa could be computed at all; the single available annotator's grades were
still used for the consensus grade in these cases, but they are excluded from the agreement
statistics above rather than silently treated as perfect agreement.

One limitation of this agreement analysis should be stated plainly: it measures agreement
\emph{between two LLM annotators}, not agreement between an LLM annotator and a human judge. It is
therefore an inter-model consistency check, evidence that the two annotators are not grading
idiosyncratically relative to each other, rather than direct evidence that either annotator's
grades track genuine human relevance judgments. This thesis does not include a human-annotated
validation sample to anchor the LLM grades against; doing so would strengthen the ground truth's
claim to validity beyond inter-annotator consistency.

\subsection{Grading the Hardened Pools}
\label{sec:methodology-groundtruth-distractors}

The native ground truth is not re-generated for the hardened pools of
Section~\ref{sec:methodology-hardened-pool}. Instead, grading proceeds incrementally: every native
argument's grade is carried over unchanged from Section~\ref{sec:methodology-groundtruth-protocol},
and only the injected distractor arguments, which by definition were never judged against the
target policy, are graded fresh, by the identical dual-annotator, MIN-consensus protocol used for
native arguments. Reusing native grades verbatim rather than re-running the annotators over them
is deliberate: because LLM annotators are not perfectly deterministic across calls, re-grading the
native arguments would perturb their existing labels for no benefit and would break the
zero-distractor condition's status as a byte-identical control
(Section~\ref{sec:methodology-hardened-levels}).

A distractor argument is graded on its merits rather than assumed irrelevant by construction.
Although a distractor is, by design, an argument that does not belong to the target policy
(Section~\ref{sec:methodology-hardened-pipeline}), a sibling-policy argument drawn for its
topical closeness to the target is not guaranteed to be irrelevant to it: stamping every
distractor grade 0 outright would not measure the hardened pool's actual relevance composition, it
would simply assert it. Grading distractors independently is therefore what makes it possible to
report, rather than assume, what fraction of an injected distractor pool a system could
legitimately have retrieved.

What the L100 grading report does establish directly is whether distractor injection achieved its
purpose: making the evaluation pool harder, in the specific sense of lowering its relevance
density below the 57.7\% figure that motivated hardening in the first place
(Section~\ref{sec:methodology-hardened-motivation}). Computed over the stratified 30-policy sample
at $D=100$, the merged native-plus-distractor pool's realized relevance base rate, 6,835 relevant
arguments (grade $\geq 2$) out of 17,820 graded pool entries, is 38.4\%, a substantial drop from
the native pool's 57.7\%. This is not uniform across policies: per-policy base rates in this
sample range from 28.5\% (``We should prohibit women in combat'') to 49.4\% (``Homeschooling
should be banned''), but every one of the 30 sampled policies falls below the native pool's
57.7\% figure, confirming that the hardening construction lowers relevance density consistently
across the sample rather than only on average.